% (c) Nikita Lisitsa, lisyarus@gmail.com, 2025

\documentclass[10pt]{beamer}

\usepackage[T2A]{fontenc}
\usepackage[russian]{babel}
\usepackage{minted}

\usepackage[most]{tcolorbox}
\tcbuselibrary{minted}

\usepackage{graphicx}
\graphicspath{ {./images/} }

\usetheme{metropolis}

\usepackage{adjustbox}

\usepackage{color}
\usepackage{soul}

\usepackage{hyperref}

\definecolor{codebg}{RGB}{29,35,49}

\setminted{fontsize=\footnotesize}

\setlength\partopsep{-\topsep}

\definecolor{blue}{rgb}{0,0,1}
\definecolor{red}{rgb}{1,0,0}

\usemintedstyle{tango}

\makeatletter
\newcommand{\slideimage}[1]{
  \begin{figure}
    \begin{adjustbox}{width=\textwidth, totalheight=\textheight-2\baselineskip-2\baselineskip,keepaspectratio}
      \includegraphics{#1}
    \end{adjustbox}
  \end{figure}
}
\makeatother

\title{Компьютерная графика}
\subtitle{Практика 4: Индексированный рендеринг, перспективная проекция, буфер глубины}
\date{2026}

\setbeamertemplate{footline}[frame number]

\begin{document}

\frame{\titlepage}

\begin{frame}[fragile]
\frametitle{Формат OBJ}
\begin{itemize}
\item Простой текстовый формат описания 3D моделей
\item Обычно используется для старых моделей, для простоты чтения и хранения
\item Довольно неоптимален по занимаемой памяти (потому что текстовый)
\item Поддерживает нормали, текстурные координаты, и простейшие материалы
\item В заготовке практики уже есть загрузчик OBJ файлов
\end{itemize}
\end{frame}

\begin{frame}[fragile]
\frametitle{Практика 4}
\begin{itemize}
\item В этой практике нельзя менять код шейдеров!
\item Все преобразования модели (сдвиги, повороты, и т.п.) и настройки камеры нужно выполнять через матрицы \mintinline{cpp}|model|, \mintinline{cpp}|view| и \mintinline{cpp}|projection|
\item В C++ матрицы заполняем по строкам, в Rust -- по столбцам (так устроены используемые библиотеки матриц)
\item \textbf{NB}: математическими библиотеками (\mintinline{rust}|glam| в Rust, \mintinline{cpp}|math| в C++) пользоваться пока нельзя :)
\end{itemize}
\end{frame}

\begin{frame}[fragile]
\frametitle{Практика 4}
\slideimage{0.png}
\end{frame}

\begin{frame}[fragile]
\frametitle{Задание 1}
Рисуем зайца
\begin{itemize}
\item В объекте \mintinline{cpp}|bunny| уже загружены вершины и индексы 3D модели
\item Настраиваем атрибуты вершин в \mintinline{cpp}|RenderPipelineDescriptor|'е (см. поля структуры \mintinline{cpp}|ObjVertex| -- нам нужны только первые два)
\item Создаём буферы для вершин (\mintinline{cpp}|usage = Vertex|) и индексов (\mintinline{cpp}|usage = Index|), загружаем в них данные (один раз, на старте программы)
\item При рисовании устанавливаем \mintinline{cpp}|render_pass|'у наши буферы: \mintinline{cpp}|render_pass.set_vertex_buffer()|, \mintinline{cpp}|render_pass.set_index_buffer()|
\item Рисуем при помощи \mintinline{cpp}|render_pass.draw_indexed()|
\end{itemize}
\end{frame}

\begin{frame}[fragile]
\frametitle{Задание 1}
\slideimage{1.png}
\end{frame}

\begin{frame}[fragile]
\frametitle{Задание 2}
Вращаем зайца
\begin{itemize}
\item Меняем матрицу \mintinline{cpp}|model| чтобы заяц крутился в \textbf{плоскости XZ}
\item В качестве угла нужно взять что-нибудь зависящее от времени, например \mintinline{cpp}|angle = time;|
\item Заяц будет обрезаться по \begin{math}z \in [0, 1]\end{math} (особенно хорошо видно, что обрезаются уши и часть головы)
\item Отмасштабируем его \textit{по всем осям} \mintinline{cpp}|scale = 0.5f|, тоже с помощью матрицы \mintinline{cpp}|model|
\end{itemize}
\end{frame}

\begin{frame}[fragile]
\frametitle{Задание 2}
\slideimage{2.png}
\end{frame}

\begin{frame}[fragile]
\frametitle{Задание 3}
Добавляем перспективу
\begin{itemize}
\item Выбираем значения \mintinline{cpp}|near|, \mintinline{cpp}|far|, \mintinline{cpp}|top|, \mintinline{cpp}|right|
\begin{itemize}
\item \mintinline{cpp}|near| -- маленький, но не слишком, в духе \begin{math}0.001 \dots 0.1\end{math}
\item \mintinline{cpp}|far| -- большой, но не слишком, в духе \begin{math}10.0 \dots 1000.0\end{math}
\item Отношение \mintinline{cpp}|right/near| -- тангенс половины угла обзора, например \mintinline{cpp}|right = near| соответствует углу обзора в \begin{math}90^\circ\end{math}
\item Отношение \mintinline{cpp}|right/top| -- aspect ratio экрана (\mintinline{cpp}|width/height|)
\end{itemize}
\item В матрицу \mintinline{cpp}|projection| записываем матрицу проекции с использованием выбранных значений (см. слайд с лекции)
\item Заяц будет виден изнутри
\end{itemize}
\end{frame}

\begin{frame}[fragile]
\frametitle{Задание 3}
\slideimage{3.png}
\end{frame}

\begin{frame}[fragile]
\frametitle{Задание 4}
Включаем тест глубины
\begin{itemize}
\item Сдвигаем камеру по оси Z на какое-то расстояние (например, на 2-3 единицы), меняя матрицу \mintinline{cpp}|view|
\begin{itemize}
\item N.B.: Камера смотрит в сторону -Z, поэтому \textit{сдвиг камеры} должен быть в сторону \textit{положительной} оси Z
\item N.B.: При этом матрица \mintinline{cpp}|view| содержит \textit{обратное преобразование!}
\end{itemize}
\item Заяц будет рисоваться неправильно: задние грани перекрывают передние
\item Настраиваем тест глубины в \mintinline{cpp}|RenderPipelineDescriptor|: \mintinline{cpp}|format = Depth24Plus|, \mintinline{cpp}|depthWriteEnabled = true| (в C++ \mintinline{cpp}|WGPUOptionalBool_True|), \mintinline{cpp}|depthCompare = Less|
\end{itemize}
\end{frame}

\begin{frame}[fragile]
\frametitle{Задание 4}
Создаём буфер глубины

% \textit{Cтрашное задание, но нужно просто много перепечатать со слайда :)}
\begin{itemize}
\item Заводим объекты \mintinline{cpp}|depthBuffer: WGPUTexture| и \mintinline{cpp}|depthBufferView: WGPUTextureView|
\item Их нужно создать на старте + пересоздавать каждый раз, когда изменились размеры экрана
\item \mintinline{cpp}|depthBuffer| создаём через \mintinline{cpp}|device.create_texture|, параметры: \mintinline{cpp}|usage = RenderAttachment|, \mintinline{cpp}|dimension = 2D|, \mintinline{cpp}|size = {width, height, 1}|, \mintinline{cpp}|format = Depth24Plus|
\item \mintinline{cpp}|depthBufferView| создаём через \mintinline{cpp}|depthBuffer.create_view|, параметры: \mintinline{cpp}{usage = RenderAttachment}, \mintinline{cpp}|dimension = 2D|, \mintinline{cpp}|format = Depth24Plus|, \mintinline{cpp}|aspect = DepthOnly|, \mintinline{cpp}|mipLevelCount = 1|, \mintinline{cpp}|arrayLayerCount = 1|
\item Добавляем \mintinline{cpp}|depthStencilAttachment| в наш \mintinline{cpp}|renderPass|: \mintinline{cpp}|view = depthBufferView|, \mintinline{cpp}|depthLoadOp = Clear|, \mintinline{cpp}|depthStoreOp = Discard|, \mintinline{cpp}|depthClearValue = 1.0|, \mintinline{cpp}|depthReadOnly = false|
\end{itemize}
\end{frame}

\begin{frame}[fragile]
\frametitle{Задание 4}
\slideimage{4.png}
\end{frame}

\begin{frame}[fragile]
\frametitle{Задание 5}
Двигаем зайца
\begin{itemize}
\item Заведём переменные \mintinline{cpp}|bunny_x| и \mintinline{cpp}|bunny_y| с координатами центра зайца по X и Y
\item Изменим матрицу \mintinline{cpp}|model|, добавив соответствующие сдвиги по X и Y
\item Перед рисованием каждого кадра обновим положение зайца:
\begin{itemize}
\item Если нажата клавиша влево (\mintinline{cpp}|SDLK_LEFT| в C++ / \mintinline{rust}|KeyCode::ArrowLeft| в Rust), сдвинем зайца влево: \mintinline{cpp}|bunny_x -= speed * dt|
\item Аналогично \mintinline{cpp}|RIGHT|, \mintinline{cpp}|DOWN|, \mintinline{cpp}|UP|
\item Состояние нажатости клавиш доступно в словаре \mintinline{cpp}|keydown| / \mintinline{cpp}|app.keydown|
\end{itemize}
\end{itemize}
\end{frame}

\begin{frame}[fragile]
\frametitle{Задание 5}
\slideimage{5.png}
\end{frame}

\begin{frame}[fragile]
\frametitle{Задание 6}
Играем с face culling
\begin{itemize}
\item Включим back-face culling: установим \mintinline{cpp}|renderPipelineDescriptor.primitive.cullMode = Back|
\item Ничего не изменится: заяц сделан так, чтобы все треугольники были CCW
\item Изменим режим на \mintinline{cpp}|Front|
\item Должны быть видны задние грани зайца и не видны передние
\item Выглядеть будет странно -- не пугайтесь, попробуйте увидеть и понять, что происходит :)
\end{itemize}
\end{frame}

\begin{frame}[fragile]
\frametitle{Задание 6}
\slideimage{6.png}
\end{frame}

\begin{frame}[fragile]
\frametitle{Задание 7*}
Три вращающихся зайца
\begin{itemize}
\item Рисуем три зайца в разных местах, вращающихся в плоскостях XY, XZ, YZ соответственно
\item Три раза настраиваем матрицу \mintinline{cpp}|model| + делаем \mintinline{cpp}|render_pass.set_immediates()| + рисуем \mintinline{cpp}|render_pass.draw_indexed()|
\item Всё ещё можно двигать стрелочками, т.е. должны учитываться \verb|bunny_x| и \verb|bunny_y| (для одного зайца, или для всех -- как хотите)
\end{itemize}
\end{frame}

\begin{frame}[fragile]
\frametitle{Задание 7*}
\slideimage{7.png}
\end{frame}


\end{document}
